\chapter[Estudos Relacionados]{Estudos Relacionados} \label{capitulo2}

Este capítulo reúne trabalhos que desenvolveram, avaliaram ou analisaram sistemas digitais destinados ao registro e à aplicação de avaliações em psicologia, saúde mental, desenvolvimento infantil e psicopedagogia, bem como sistemas web voltados à educação. Após a descrição do procedimento de seleção (Seção~\ref{sec:selecao_estudos}), os estudos são discutidos em dois grupos, conforme sua proximidade com o objeto desta pesquisa (Seções~\ref{sec:estudos_avaliacao} e~\ref{sec:estudos_educacao}). Em seguida, apresenta-se um levantamento das soluções comerciais disponíveis no mercado brasileiro (Seção~\ref{sec:solucoes_mercado}) e, por fim, uma síntese que delimita a lacuna abordada neste trabalho (Seção~\ref{sec:sintese_comparativa}).

\section{Seleção dos estudos}\label{sec:selecao_estudos}

A seleção dos estudos ocorreu em duas etapas. Na primeira, de caráter exploratório, foram reunidos trabalhos sobre sistemas web aplicados à gestão escolar e à prática docente, os quais forneceram subsídios para requisitos gerais de controle de acesso, organização de registros e usabilidade. Como esses trabalhos não abordavam sistemas de avaliação, realizou-se uma segunda etapa, orientada a estudos mais próximos do objeto desta pesquisa.

Nessa etapa, as buscas foram conduzidas nas bases SciELO, PePSIC e Biblioteca Digital Brasileira de Teses e Dissertações (BDTD) e complementadas por buscas abertas na web, a fim de alcançar a literatura internacional e trabalhos depositados em repositórios institucionais. Empregaram-se, em português e em inglês, as seguintes combinações de termos: ``prontuário eletrônico'' e ``psicologia''; ``avaliação psicopedagógica'' associada a ``\textit{software}'', ``sistema'', ``digital'' ou ``informatizado''; ``triagem do desenvolvimento infantil'' associada a ``aplicativo'', ``digital'', ``eletrônico'' ou ``\textit{online}''; ``\textit{Ages and Stages Questionnaire}'' associado a ``\textit{web-based}'' ou ``\textit{electronic}''; e ``\textit{electronic health records}'' associado a ``\textit{mental health}''.

Foram considerados trabalhos publicados entre 2013 e 2025, em português, inglês ou espanhol, que descrevessem o desenvolvimento ou a avaliação de sistemas digitais para registro ou aplicação de avaliações nas áreas mencionadas, ou que revisassem a literatura sobre sua implementação. Excluíram-se jogos educacionais, \textit{softwares} de intervenção e demais recursos sem relação com o registro do processo avaliativo. Ressalta-se que a busca não teve caráter de revisão sistemática, uma vez que seu propósito foi identificar trabalhos representativos de cada linha de pesquisa, e não esgotar a produção sobre o tema.

\section{Sistemas de registro e de avaliação}\label{sec:estudos_avaliacao}

Os estudos que tratam diretamente do registro ou da aplicação de avaliações podem ser organizados em três linhas: os prontuários eletrônicos em psicologia e saúde mental, as ferramentas digitais de triagem do desenvolvimento infantil e o uso de \textit{software} na avaliação psicopedagógica.

No âmbito dos prontuários eletrônicos, \citeonline{silva2020_sistema_psi} desenvolveram o Sistema PSI, destinado a psicólogos que atuam em mais de um local de trabalho e que permite o cadastro de clínicas, equipes, pacientes e sessões. Embora o sistema tenha sido testado e aprovado por profissionais da área, o prontuário é concebido como registro de sessões, sem que a aplicação de instrumentos avaliativos seja estruturada. A dissertação de \citeonline{valgas2021_prontuario_neuropsicologico} constitui, entre os trabalhos encontrados, o mais próximo desta pesquisa. Fundamentado igualmente na \textit{Design Science Research}, o autor propôs um sistema de gestão da avaliação neuropsicológica em contexto hospitalar e identificou, com os profissionais, as tarefas mais onerosas do processo, entre elas a pontuação de testes, a alternância entre diferentes sistemas e a transcrição de informações registradas em papel. O protótipo obteve avaliação de usabilidade positiva, e o estudo conclui pela necessidade de padrões de interoperabilidade para os dados neuropsicológicos. Distingue-se do presente trabalho, contudo, pelo público atendido, composto por adultos hospitalizados, e pelo campo profissional, a neuropsicologia.

Duas revisões de escopo contribuem para compreender os desafios da adoção desses sistemas em saúde mental. Ao analisarem 23 estudos sobre a implantação de prontuários eletrônicos nesses serviços, \citeonline{zurynski2021_emr_saude_mental} verificaram que a adoção depende de fatores financeiros, técnicos e organizacionais, mas sobretudo da percepção dos profissionais de que o sistema não interfere negativamente em seu fluxo de trabalho, razão pela qual recomendam que tais sistemas sejam concebidos com a participação de seus usuários. \citeonline{kariotis2022_ehr_saude_mental}, por sua vez, evidenciaram que a exigência de registro em campos estruturados entra em tensão com o caráter narrativo das informações em saúde mental, de modo que detalhes contextuais tendem a se perder e informações sensíveis acabam omitidas ou registradas fora do sistema.

Na triagem do desenvolvimento infantil, \citeonline{moreira2022_tedi} apresentaram o TEDI, ferramenta eletrônica brasileira de acesso livre, elaborada com a participação de especialistas e submetida a testes de usabilidade e de confiabilidade. Em âmbito internacional, \citeonline{schonhaut2023_asq_web} avaliaram a aplicação \textit{web} da versão em espanhol do ASQ-3 no Chile, a partir de 2.318 questionários referentes a 1.061 crianças, e constataram que o formato é confiável e bem aceito pelos pais. Já no contexto escolar, \citeonline{Ramos_DellaBarba_2021} demonstraram que professoras de centros municipais de educação infantil foram capazes de aplicar a versão brasileira do ASQ a 392 crianças. A revisão de \citeonline{pereira2021_rastreio_movel}, por outro lado, identificou apenas seis estudos sobre aplicativos de vigilância e rastreio do desenvolvimento, realizados majoritariamente em serviços de saúde, e destacou a escassez de investigações sobre o impacto de seu uso. Observa-se que os trabalhos dessa linha digitalizam um instrumento específico e previamente validado, com itens e critérios de pontuação fixos, sem oferecer ao profissional a possibilidade de cadastrar ou adaptar instrumentos e sem integrar os resultados a registros qualitativos e relatórios interpretativos.

Quanto ao uso de \textit{software} na avaliação psicopedagógica, apenas um trabalho foi localizado. \citeonline{pinheiro2013_avaliacao_psicopedagogica} investigou a aplicabilidade do \textit{software} Desenvolve\textsuperscript{\textregistered} na avaliação de 16 crianças de cinco a sete anos em acolhimento institucional. O programa afere o desempenho em 19 habilidades cognitivas e apresenta os resultados em gráficos, mas caracteriza-se como instrumento de avaliação, e não como sistema de registro e acompanhamento do processo avaliativo. Cabe notar que a própria autora registrou a carência de documentação da avaliação pedagógica em prontuário, aspecto que se aproxima diretamente da lacuna explorada neste trabalho.

\section{Sistemas web aplicados à educação}\label{sec:estudos_educacao}

Os trabalhos reunidos na primeira etapa da busca contribuem sobretudo com requisitos de caráter geral. \citeonline{abreu2025_sistemas_web_gestao} e \citeonline{revistaft2024_sistema_vida_escolar} descrevem sistemas de gestão escolar e apontam o controle de acesso, a padronização dos registros e a capacitação dos usuários como condições para a efetiva adoção dessas ferramentas. \citeonline{pedroso2025_integraedu} apresentam o IntegraEDU, plataforma de análise de dados educacionais que oferece gestão de perfis e relatórios individualizados, enquanto \citeonline{carneiro2023_taep4} documentam a plataforma TAEP4.0, voltada ao registro das práticas docentes. O aplicativo desenvolvido por \citeonline{silva2024_app_tdah} para apoiar professores na identificação de indicativos de TDAH contribui, principalmente, no aspecto metodológico, por ter sido construído por meio de prototipação iterativa e validado com especialistas. Por fim, \citeonline{Aielo_Silva_Augusto_2025} discutem, em perspectiva teórica, o potencial das tecnologias digitais para o acompanhamento individualizado na intervenção psicopedagógica. Embora relevantes para o projeto da interface e do controle de acesso, nenhum desses trabalhos contempla o registro de avaliações psicopedagógicas.

\section{Soluções existentes no mercado}\label{sec:solucoes_mercado}

Com o intuito de verificar se a lacuna observada na literatura se mantém entre as soluções comerciais, realizou-se um levantamento exploratório de sistemas de prontuário eletrônico oferecidos no mercado brasileiro a psicopedagogos, psicólogos e clínicas multiprofissionais. A análise baseou-se nas funcionalidades divulgadas nos sítios oficiais de cada produto, consultados em novembro de 2026. Por esse motivo, a ausência de um recurso no Quadro~\ref{qua:solucoes_mercado} indica apenas que ele não é informado pelo fornecedor, e não necessariamente que inexista no produto.

O iPsy Pro, destinado a psicopedagogos, reúne agenda, controle financeiro e um prontuário estruturado em resumo, anamnese, avaliação, plano e sessões, além de modelos de documentos e de um banco de jogos para uso em atendimento. O próprio fornecedor esclarece que a ferramenta não substitui a avaliação nem os instrumentos de uso restrito, o que indica a inexistência de aplicação estruturada de protocolos no sistema \cite{ipsy_pro}. O Cerebraris, voltado a psicopedagogos e neuropsicopedagogos, avança nesse aspecto ao permitir a aplicação de um conjunto predefinido de instrumentos, como o TDE-II, o SNAP-IV e o M-CHAT, com pontuação automática e emissão de laudos a partir de modelos \cite{cerebraris}; não há, entretanto, informação sobre o cadastro de instrumentos pelo próprio profissional. Entre as soluções de caráter generalista, o PsicoManager integra prontuário, escalas, agenda, financeiro e permissões por usuário \cite{psicomanager}, e o Ninsaúde Clinic permite a personalização de formulários de prontuário a partir de modelos criados pelo usuário \cite{ninsaude_prontuario}. Nesse último caso, contudo, não se informa de que modo a alteração de um modelo repercute nos registros já preenchidos.

\begin{quadro}[H]
\centering
\footnotesize
\caption{Funcionalidades divulgadas por soluções de prontuário eletrônico no mercado brasileiro}
\begin{tabular}{|>{\raggedright\arraybackslash}p{2.0cm}|>{\raggedright\arraybackslash}p{2.5cm}|>{\raggedright\arraybackslash}p{2.6cm}|>{\raggedright\arraybackslash}p{2.2cm}|>{\raggedright\arraybackslash}p{2.4cm}|>{\raggedright\arraybackslash}p{1.6cm}|}
\hline
\textbf{Produto} & \textbf{Público} & \textbf{Instrumentos de avaliação} & \textbf{Criação de instrumentos próprios} & \textbf{Versão vinculada às aplicações} & \textbf{Agenda e financeiro} \\
\hline
iPsy Pro & Psicopedagogos & Não (declarado pelo fornecedor) & Não informado & Não informado & Sim \\ \hline
Cerebraris & Psicopedagogos e neuropsico\-pedagogos & Conjunto predefinido, com pontuação automática & Não informado & Não informado & Sim \\ \hline
PsicoManager & Psicólogos e clínicas & Escalas (sem detalhamento) & Não informado & Não informado & Sim \\ \hline
Ninsaúde Clinic & Clínicas multiprofissionais & Formulários de prontuário & Sim (modelos de formulário) & Não informado & Sim \\ \hline
Sistema proposto & Psicopedagogos & Cadastrados pelo profissional & Sim & Sim & Não \\
\hline
\end{tabular}
\fonte{Elaborado pela autora (2026), com base em \citeonline{ipsy_pro}, \citeonline{cerebraris}, \citeonline{psicomanager} e \citeonline{ninsaude_prontuario}.}
\label{qua:solucoes_mercado}
\end{quadro}

Depreende-se do levantamento que já existem sistemas específicos para a psicopedagogia e que estes atendem satisfatoriamente às rotinas administrativas do consultório, como agendamento, cobrança e emissão de documentos, funcionalidades que não integram o escopo deste trabalho. A lacuna identificada é, portanto, mais restrita e refere-se ao tratamento do instrumento avaliativo. Nenhuma das soluções analisadas divulga, simultaneamente, a possibilidade de o profissional cadastrar e adaptar seus próprios instrumentos e um controle de versões que preserve, em cada aplicação, a versão efetivamente utilizada, condição necessária para comparar aplicações realizadas em momentos distintos com a certeza de quais questões foram respondidas em cada uma.

\section{Síntese e delimitação da lacuna}\label{sec:sintese_comparativa}

O Quadro~\ref{qua:estudos_relacionados} sintetiza os trabalhos acadêmicos discutidos neste capítulo, indicando a principal contribuição de cada um e o aspecto em que se diferenciam da proposta aqui apresentada.

{\footnotesize
\setlength{\LTpre}{0pt}
\setlength{\LTpost}{0pt}
\begin{center}\captionsetup{hypcap=false}\captionof{quadro}{Síntese comparativa dos estudos relacionados}\label{qua:estudos_relacionados}\end{center}
\vspace{-1.2\baselineskip}
\begin{longtable}{|>{\raggedright\arraybackslash}p{2.6cm}|>{\raggedright\arraybackslash}p{2.8cm}|>{\raggedright\arraybackslash}p{4.2cm}|>{\raggedright\arraybackslash}p{4.2cm}|}
\hline
\textbf{Autor e ano} & \textbf{Tipo de trabalho} & \textbf{Contribuição principal} & \textbf{Diferença em relação a este trabalho} \\ \hline
\endfirsthead
\multicolumn{4}{r}{\textit{(continuação)}} \\ \hline
\textbf{Autor e ano} & \textbf{Tipo de trabalho} & \textbf{Contribuição principal} & \textbf{Diferença em relação a este trabalho} \\ \hline
\endhead
\multicolumn{4}{r}{\textit{(continua)}} \\
\endfoot
\endlastfoot
Silva et al. (2020) & Prontuário \textit{web} para psicólogos & Digitalização do registro de sessões & Não estrutura a aplicação de instrumentos \\ \hline
Valgas (2021) & Dissertação; sistema de avaliação neuropsicológica & Gestão da avaliação com \textit{Design Science Research} e avaliação de usabilidade & Público adulto hospitalar; campo da neuropsicologia \\ \hline
Zurynski et al. (2021) & Revisão de escopo & Fatores de adoção de prontuários em saúde mental & Não propõe sistema \\ \hline
Kariotis et al. (2022) & Revisão de escopo & Tensão entre campos estruturados e registro narrativo & Não propõe sistema \\ \hline
Moreira et al. (2022) & Ferramenta de triagem (TEDI) & Triagem do desenvolvimento infantil em meio eletrônico & Instrumento fixo; sem registros qualitativos \\ \hline
Schonhaut et al. (2023) & Estudo de validação & Confiabilidade do ASQ-3 aplicado pela \textit{web} & Instrumento fixo; aplicação pelos pais \\ \hline
Ramos e Della Barba (2021) & Estudo de aplicação & Uso do ASQ-BR por professoras da educação infantil & Aplicação em papel; não propõe sistema \\ \hline
Pereira, Charchat-Fichman e Landeira-Fernandez (2021) & Revisão & Panorama de aplicativos de rastreio do desenvolvimento & Não propõe sistema \\ \hline
Pinheiro (2013) & Dissertação; \textit{software} na avaliação psicopedagógica & Uso de \textit{software} na avaliação; aponta falta de registro em prontuário & \textit{Software} de avaliação, não de registro \\ \hline
Abreu et al. (2025) & Sistema de gestão escolar & Controle de acesso e padronização de registros & Não aborda avaliação \\ \hline
Sousa, Sousa e Sousa (2024) & Sistema de gestão escolar & Lições de usabilidade e capacitação & Não aborda avaliação individualizada \\ \hline
Pedroso et al. (2025) & Plataforma de dados educacionais & Perfis e relatórios individualizados & Foco administrativo \\ \hline
Carneiro et al. (2023) & Plataforma docente & Registro estruturado de práticas & Não trata de dados clínico-pedagógicos \\ \hline
Magalhães et al. (2024) & Aplicativo para professores & Prototipação e validação com especialistas & Foco restrito ao TDAH \\ \hline
Aielo, Silva e Augusto (2025) & Estudo teórico & Tecnologia na intervenção psicopedagógica & Não apresenta sistema \\ \hline
\end{longtable}
\vspace{-0.5\baselineskip}
\fonte{Elaborado pela autora (2026).}
}

A análise dos estudos revela que a digitalização do registro clínico e da triagem do desenvolvimento infantil já constitui objeto de pesquisa, com resultados favoráveis quanto à confiabilidade, à usabilidade e à aceitação. Os prontuários eletrônicos em psicologia, entretanto, tratam o registro de modo genérico, ao passo que as ferramentas de triagem digitalizam instrumentos fixos. O trabalho mais próximo, de \citeonline{valgas2021_prontuario_neuropsicologico}, dirige-se à neuropsicologia de adultos, e o único estudo localizado sobre \textit{software} na avaliação psicopedagógica \cite{pinheiro2013_avaliacao_psicopedagogica} emprega um instrumento de avaliação e aponta, ele próprio, a falta de registro do processo. As soluções comerciais, por sua vez, atendem à gestão do consultório, mas não divulgam o cadastro de instrumentos acompanhado de controle de versões vinculado às aplicações.

Delimita-se, assim, a lacuna que orienta este trabalho: a inexistência, entre os estudos e produtos analisados, de um sistema que reúna, no contexto da avaliação psicopedagógica na Educação Infantil, o cadastro de instrumentos pelo próprio profissional, o versionamento capaz de preservar a versão utilizada em cada aplicação, o registro das respostas e a elaboração de relatórios qualitativos. A literatura examinada oferece, ainda, duas orientações para o projeto. A primeira diz respeito à participação dos profissionais na concepção do sistema \cite{zurynski2021_emr_saude_mental}, buscada neste trabalho por meio da entrevista de levantamento de requisitos e do questionário de avaliação. A segunda refere-se ao equilíbrio entre campos estruturados e espaços de registro narrativo \cite{kariotis2022_ehr_saude_mental}, que se reflete na combinação entre as respostas aos instrumentos e os relatórios redigidos em texto livre.